\documentclass[../sparc.tex]{subfiles}
\graphicspath{{\subfix{../images/}}}
\begin{document}

\newpage
%%%%%%%%%%%%%%%%%%%%%%%%%%%%%%%%%%%%%%%%%%%%%%%%%%%%%%%%%%%%%%%%%%%%%%%%%%%%%%%%
\section{Serial Peripheral Interface (SPI) bus}
\label{section:spi}
\index{Electronics!SPI Bus}
\newglossaryentry{SPI}{name=SPI, description={Serial Peripheral Interface}}

%%%%%%%%%%%%%%%%%%%%%%%%%%%%%%%%%%%%%%%%%%%%%%%%%%%%%%%%%%%%%%%%%%%%%%%%%%%%%%%%
\subsection{General information}

\textit{\gls{SPI}} -- a synchronous data transfer bus that is used for
communication between devices on short distances.

The bus uses ``controller-periphery'' model where the controller controls the
communication with a periphery device, initiating the data transfer.  There can
be several periphery devices.  Some devices can change the role between
``controller'' and ``periphery'' on the fly.

The protocol works in full duplex mode and uses four \cite{arduino:spi} logical
lines:

\begin{itemize}
\item \textbf{CS} (Chip Select) -- the signal from the controller that initiates
  the communication with a selected periphery device on the SPI bus.  The line
  is active when its value is low.  There must be one CS line for each periphery
  device, thus the number of such lines depends on the number of periphery
  devices.
\item \textbf{SCLK} (Serial Clock) -- the clock signal from the controller.
\item \textbf{COPI} (Controller Out Peripheral In) -- the line that is used to
  transfer data from the controller to a periphery device.
\item \textbf{CIPO} (Controller In, Peripheral Out) -- the line that is used to
  transfer data from a periphery device to the controller.
\end{itemize}

In the past, terms ``master'' and ``slave'' were used in the place of
``controller'' and ``periphery''.  In this book, we will be using modern
terminology ``controller''/''periphery'' -- the table
\ref{table:spi-old-terminology} shows the correspondence between the old terms
of SPI bus with the new ones.

\begin{table}[H]
  \centering
  \begin{tabular}{ | c | c | }
    \hline
    \textbf{Master/Slave (old variant)}
    & \textbf{Controller/Periphery (new variant)} \\
    \hline
    Master In Slave Out (MISO) & Controller In, Peripheral Out (CIPO) \\
    \hline
    Master Out Slave In (MOSI) & Controller Out Peripheral In (COPI) \\
    \hline
    Slave Select pin (SS) & Chip Select Pin (CS) \\
    \hline
  \end{tabular}
  \caption{The correspondence of the old SPI terms with the new ones.}
  \label{table:spi-old-terminology}
\end{table}

The SPI bus can be schematically represented as is shown on the
fig. \ref{fig:spi-bus}.

\figureSPI{en}

The communication between the controller and a periphery device can be
represented as a system where both sides have a data buffer and the data
transfer occurs between those two buffers.

There can be differences in digital ports assigned for the
SPI\cite{arduino:reference} protocol on different the Arduino platforms, as is
shown in the table \ref{table:spi-ports}.

\tableSPIPorts{en}

It should be noted that the implementation of the SPI bus can vary, thus when
working with each new device through the SPI bus we have to rely on the
documentation for that device.

%%%%%%%%%%%%%%%%%%%%%%%%%%%%%%%%%%%%%%%%%%%%%%%%%%%%%%%%%%%%%%%%%%%%%%%%%%%%%%%%
\subsection{SPI library}

%%%%%%%%%%%%%%%%%%%%%%%%%%%%%%%%%%%%%%%%%%%%%%%%%%%%%%%%%%%%%%%%%%%%%%%%%%%%%%%%
\subsubsection{Including the library to a project}

The Arduino IDE ships with the SPI library\cite{arduino:reference}, that
provides procedures for working with the bus.  We can add the library into a
project by including the ``spi.h'' header:

\begin{minted}{cpp}
  #include <SPI.h>
\end{minted}

%%%%%%%%%%%%%%%%%%%%%%%%%%%%%%%%%%%%%%%%%%%%%%%%%%%%%%%%%%%%%%%%%%%%%%%%%%%%%%%%
\subsubsection{Configuration of SPI parameters}

The configuration of the SPI bus is performed through the
\mintinline{cpp}{SPISettings} class which allows us to set three options:

\begin{itemize}
\item Maximum SPI speed.  The value of this parameter depends on the clock
  frequency of the CPU -- for example, when the frequency is 15MHz, the SPI
  frequency must be set to 15'000'000.  Arduino automatically set the optimal
  bus speed which equals to or less than the specified frequency.
\item The bit order.  There are two variants: the most significant bit (MSB)
  first, or the least significant bit (LSB) first.  In the SPI options those
  variants can be set with \mintinline{cpp}{MSBFIRST} and
  \mintinline{cpp}{LSBFIRST} constants respectively.
\item The SPI working mode.  There are four modes that control whether the bits
  are shifted on rising or falling edges of clock pulses, and whether the idle
  value of the clock line is \texttt{LOW} or \texttt{HIGH}.
\end{itemize}

\begin{table}[H]
  \small
  \centering
  \begin{tabular}{ | m{1.75cm} | m{1cm} | m{1cm} | m{1cm} | m{1cm} | }
    \hline
    \textbf{Режим}
    & \textbf{Clock Polarity (CPOL)}
    & \textbf{Clock Phase (CPHA)}
    & \textbf{Output Edge}
    & \textbf{Data Capture} \\
    \hline
    \texttt{SPI\_MODE0} & 0 & 0 & Falling & Rising \\
    \hline
    \texttt{SPI\_MODE1} & 0 & 1 & Rising & Falling \\
    \hline
    \texttt{SPI\_MODE2} & 1 & 0 & Rising & Falling \\
    \hline
    \texttt{SPI\_MODE3} & 1 & 1 & Falling & Rising \\
    \hline
  \end{tabular}
  \caption{SPI working modes.}
  \label{table:spi-modes}
\end{table}

The all four SPI working modes are shown in the table \ref{table:spi-modes}.
All the mentioned parameters are set in the \mintinline{cpp}{SPISettings} class
constructor at the time of the instance creation.  For example, if we want to
use the SPI frequency bus equal to 14'000'000 Hz, the MSB bit order and the SPI
mode set to \mintinline{cpp}{SPI_MODE0}, we can create
\mintinline{cpp}{SPISettings} instance as follows:

\begin{listing}[H]
  \begin{minted}{cpp}
    SPISettings settings(14000000, MSBFIRST, SPI_MODE0);
  \end{minted}
  \label{listing:communication-spi-settings}
  \caption{Creation of \mintinline{cpp}{SPISettings} with the required
    parameters.}
\end{listing}

%%%%%%%%%%%%%%%%%%%%%%%%%%%%%%%%%%%%%%%%%%%%%%%%%%%%%%%%%%%%%%%%%%%%%%%%%%%%%%%%
\subsubsection{Data transfer}

The interaction with the SPI bus is done through the global
\mintinline{cpp}{SPIClass} instance which is called \mintinline{cpp}{SPI}.  This
object become available as soon as we include \mintinline{cpp}{spi.h} header
file into our project.

The configuration of the SPI bus starts with the calling to
\mintinline{cpp}{begin} method from the \mintinline{cpp}{SPI} instance.  This
must be done once, thus the preferred place to do this is the
\mintinline{cpp}{setup} procedure.

Next we can use the created \mintinline{cpp}{SPISettings} instance (let's call
it \mintinline{cpp}{settings}) to start the SPI communication session.

To start the communication session on the SPI bus we have to acquire an
exclusive access to the bus using \mintinline{cpp}{beginTransaction} method.  We
have to pass an \mintinline{cpp}{SPISettings} instance to this method.  The call
to the method stores the SPI settings for further usage in the global SPI
instance.

After that we can transfer data through the SPI bus using the methods of data
transfer, which also available from the \mintinline{cpp}{SPI} instance.

When data transfer is done, we have to call \mintinline{cpp}{endTransaction}
method to finish the SPI communication session.  If we fail to do that, it can
disturb the workings of the other code parts and libraries that also work with
the SPI bus.

If SPI bus is used to communicate with two or more devices, we have to select
the right periphery device for communication using of ``Chip Select'' (CS)
lines, by pulling the right line to the \texttt{LOW} value.  After the
communication with the periphery device is done, we have to pull the
corresponding CS line to \texttt{HIGH}.

Let's assume that we want to transfer one byte of data (number 42) to a
periphery device, which CS line is connected to the digital port number 2.

\begin{listing}[H]
  \begin{minted}{cpp}
    const byte CS = 2;
    SPISettings settings(14000000, MSBFIRST, SPI_MODE0);
    void setup() {
      pinMode(CS, OUTPUT);

      SPI.begin();

      SPI.beginTransaction(settings);
      digitalWrite(CS, LOW);
      uint8_t return_value = SPI.transfer(42);
      digitalWrite(CS, HIGH);
      SPI.endTransaction();
    }
  \end{minted}
  \label{listing:communication-spi-begin}
  \caption{An example of working with the SPI bus.}
\end{listing}

In the example \ref{listing:communication-spi-begin} we are using SPI bus in the
\mintinline{cpp}{setup} procedure, although we can do that the same way in the
\mintinline{cpp}{loop} or any other procedure.  Using \mintinline{cpp}{transfer}
method we pass the value 42 to a periphery device, and get one byte in response.
The response is saved in \mintinline{cpp}{return_value} variable -- we can ignore
this value if we don't expect anything from the periphery device in the
response.

It is worth mentioning that there are three methods for data transfer:

\begin{itemize}
\item \mintinline{cpp}{uint8_t transfer(uint8_t data)} -- a method for
  transferring of an one-byte value to a periphery device over the COPI line.
  The method returns the value acquired from the device over the CIPO line.
\item \mintinline{cpp}{uint16_t transfer16(uint16_t data)} -- a method for
  transferring a two-byte value to a periphery device.  As for the previous
  method, the procedure returns the value acquired from the device over the CIPO
  line, but a two-byte value in this case.
\item \mintinline{cpp}{void transfer(void *buf, size_t count)} -- a method for
  transferring a data buffer \mintinline{cpp}{buf} of \mintinline{cpp}{count}
  byte in size.  The response from the device over the CIPO line is written to
  the buffer \mintinline{cpp}{buf} in the place of each transferred byte.
\end{itemize}

%%%%%%%%%%%%%%%%%%%%%%%%%%%%%%%%%%%%%%%%%%%%%%%%%%%%%%%%%%%%%%%%%%%%%%%%%%%%%%%%
\subsection{Controlling of a shift register}

As we mentioned in the \ref{section:spi-basics} section, the communication
between a controller and a periphery device over the \gls{SPI} bus can be
represented as the system of two data buffers between which the data transfer
occurs.  The operation of \gls{SIPO} and \gls{PISO} shift registers fits well
into this concept.

Let's take \gls{SIPO} for example -- it represents an one-byte (8 bits) external
memory into which we sequentially shift out the bits over the ``DS'' line,
acquiring the shifted out bits from the register over the ``QA'' line.  It can
be said that we have an one-byte memory on the controller side as well, from
which we shift out the data to the register.  If we ``loop'' those two memory
buffers together so the shifted out bits from the register are written in the
controller memory, we will effectively get the SPI bus implementation.

\figureSPIShiftRegister{en}

A schematic representation of a buffer on the controller side (``Controller'')
with the buffer on the shift register side (``Periphery'') is shown in the
fig. \ref{fig:spi-shift-register}.

\figureSPIShiftRegisterConnections{en}

The connection between a controller and a SIPO shift register is shown in the
fig. \ref{fig:spi-sipo-register}.

In addition to the ``COPI'' and ``CIPO'' lines, the controller have to be
connected to the shift register with the ``LATCH'' line, which is used to output
the data from the shift register memory to its outputs ``Q0...Q7''.  This line
replaces the standard ``chip select'' (``CS'') line which is used in the SPI
bus.

\begin{listing}[ht]
  \begin{minted}{cpp}
    #include <SPI.h>

    const byte LATCH = 4;
    SPISettings settings(2000000, MSBFIRST, SPI_MODE2);

    void setup() {
      Serial.begin(9600);
      pinMode(LATCH, OUTPUT);
      SPI.begin();
    }

    int data_index = 0;
    const char DATA[] = "hello";
    void loop() {
      char response;

      SPI.beginTransaction(settings);
      response = SPI.transfer(DATA[data_index]);
      SPI.endTransaction();
      digitalWrite(LATCH, HIGH);
      delay(1);
      digitalWrite(LATCH, LOW);
      Serial.print("character sent:     ");
      Serial.println(DATA[data_index]);
      Serial.print("character received: ");
      Serial.println(response);
      delay(500);
      data_index++;
      if (data_index == sizeof(DATA)) {
        data_index = 0;
      }
    }
  \end{minted}
  \label{listing:communication-spi-shift-register}
  \caption{An example of communication with a shift register using the SPI bus.}
\end{listing}

An example of data transfer between a shift register and a controller is shown
in the listing \ref{listing:communication-spi-shift-register}.  We are using
``hello'' string as the data sample, transferring it letter-by-letter to the
shift register.  After transferring each byte (one symbol) to the register, we
are switching ``LATCH'' line, thus outputting the binary code of the symbol in
the register memory to its output pins.  In the response, we are receiving the
previous byte (a symbol) from the register memory.  For sake of clarity, we
print the transferred and received data to the serial port.

Application of the SPI bus for the communication with shift register allows us
to avoid writing our own code for the data transfer, thus making our programs
more readable.

%%%%%%%%%%%%%%%%%%%%%%%%%%%%%%%%%%%%%%%%%%%%%%%%%%%%%%%%%%%%%%%%%%%%%%%%%%%%%%%%
\subsubsection{Exercises}

\begin{enumerate}
\item Connect a SPI shift register using SPI bus to an Arduino.  Implement
  various LED effects using the shift register.
\item Connect two shift registers to an Arduino using SPI bus.  Implement a
  binary counter which outputting numbers in the binary form to a shift register
  and to LEDs.  Sequentially output numbers from 0 to $2^{16} - 1$ to the counter.
\item Modify the previous task by adding the counter speed regulation by means
  of a potentiometer.
\item Investigate the counter overflow in the second task when using
  \mintinline{cpp}{int} data type on the Arduino.  How the number is encoded in
  the binary form after crossing the $2^{15} - 1$ threshold?
\end{enumerate}

\end{document}
